\label{sec:clifford-area-transfer}

A source-radius factor compensates for the dilation term in
\eqref{eq:pullback-defect} in sufficiently large admissible dimensions.

\begin{theorem}[Area minimization under Clifford pullback]
\label{thm:clifford-area-transfer}
Let \(f\) be a nonzero homogeneous polynomial on \(\R^m\), and let
\(Z\subset\{f=0\}\) be a closed conical real algebraic set containing
the origin and having codimension at least two.  Suppose that
\(H\in C^3(\R^m\setminus Z)\cap C^0(\R^m)\) is homogeneous of
degree \(\ell>0\), with
\[
 \operatorname{sgn}H=\operatorname{sgn}f,\qquad
 |\nabla_yH|>0\quad\text{on }\R^m\setminus Z.
\]
Put \(W_H=|\nabla_yH|^2\), and suppose that, off \(Z\),
\[
 H\cL_y(H)\geq0,\qquad H\Delta_yH\geq-MW_H
\]
for some finite \(M\geq0\).  Let \(P:\R^N\to\R^m\) be the
Clifford map above.  If \(a>0\) and
\begin{equation}\label{eq:clifford-transfer-bound}
 N\geq1+4\ell(1+M)+2a(1+M)+\frac{\ell}{a},
\end{equation}
then
\[
 G=R^aH(P),\qquad E=\{f(P)<0\},
\]
defines a two-sided subcalibration: off \(P^{-1}(Z)\),
\(\nabla_xG\ne0\) and \(G\cL_x(G)\geq0\), with strict inequality
in the two open phases.  Consequently \(E\) is locally perimeter
minimizing and \(\partial\llbracket E\rrbracket\) is locally mass
minimizing. If \(N\geq3\) and both phases are nonempty, this current
is strictly area minimizing. If both signs of \(f\) occur in every
neighborhood of each target zero, then
\[
 \operatorname{spt}\partial\llbracket E\rrbracket=\{f(P)=0\}.
\]
In particular, this last condition holds when \(f\) is harmonic.
\end{theorem}

\begin{proof}
Set \(h=H(P)\) and \(b=a(a+2\ell)\).  The chain rule and Euler's
identity give, on \(\R^N\setminus P^{-1}(Z)\),
\begin{align}
 |\nabla_xG|^2
 &=4R^{2a+1}W_H(P)+4bR^{2a-1}h^2,
 \label{eq:clifford-transfer-gradient}\\
 \cL_x(G)
 &=16R^{3a+2}(\cL_yH)(P)
   +8\bigl[a(N-1)-\ell-2b\bigr]R^{3a}hW_H(P)\notag\\
 &\quad+16bR^{3a}h^2(\Delta_yH)(P)
   +8b\bigl[a(N-1)+\ell\bigr]R^{3a-2}h^3.
 \label{eq:clifford-transfer-defect}
\end{align}
For the second formula, use the radial identity
\eqref{eq:radial-master-L} with the source function \(h\), of degree
\(2\ell\), and substitute
\[
\begin{gathered}
 |\nabla_xh|^2=4RW_H(P),\qquad
 \Delta_xh=4R(\Delta_yH)(P),\\
 \cL_x(h)=16R^2(\cL_yH)(P)-8\ell hW_H(P).
\end{gathered}
\]
These identities follow from \eqref{eq:Clifford-pullback-identities}.

Multiplication of \eqref{eq:clifford-transfer-defect} by \(h\) yields
\begin{align*}
 h\cL_x(G)\geq{}&
 8\bigl[a(N-1)-\ell-2b(1+M)\bigr]R^{3a}h^2W_H(P)\\
 &+8b\bigl[a(N-1)+\ell\bigr]R^{3a-2}h^4.
\end{align*}
The first coefficient is nonnegative by
\eqref{eq:clifford-transfer-bound}, and the last term is positive
when \(h\ne0\).  Since \(R^a>0\), this proves the claimed phase
inequality.  Formula \eqref{eq:clifford-transfer-gradient} gives
\(\nabla_xG\ne0\) off \(Z'=P^{-1}(Z)\).

The extension of \(G\) is continuous with the signs of \(f(P)\);
at the origin, use \(H(P(x))=O(|x|^{2\ell})\).
Surjectivity gives \(f(P)\not\equiv0\), hence local finite perimeter.
Since \(P\) is a submersion off the origin, stratifying \(Z\) shows
that \(Z'=P^{-1}(Z)\) has codimension at least two. Thus
\cref{cor:gradient-subcalibration} applies. For \(N\ge3\) and two
nonempty phases, the strict divergence signs give strict minimality
by \cref{thm:strict-minimality-flux}.

The additional phase assumption transfers by submersion away from
the origin and by surjectivity and homogeneity at the origin.
The support conclusion follows from \cref{rem:polynomial-finite-perimeter}.
\end{proof}

For harmonic nonconstant \(f\), \eqref{eq:Clifford-pullback-identities}
and surjectivity make \(F=f\circ P\) nonzero and harmonic.
Area minimization then implies real irreducibility by
\cref{prop:minimizing-irreducibility}. Nonharmonic targets still
require the stated phase hypothesis for full support.

The ratio bound is preserved under iteration:
\[
 \Delta_xG=4R^{a+1}(\Delta_yH)(P)
       +2a(N+2a+4\ell-2)R^{a-1}h,
\]
and \eqref{eq:clifford-transfer-gradient} gives
\begin{align}
 G\Delta_xG+M|\nabla_xG|^2
 ={}&4R^{2a+1}\bigl(H\Delta_yH+MW_H\bigr)(P)\notag\\
 &+\bigl[2a(N+2a+4\ell-2)+4Mb\bigr]R^{2a-1}h^2
 \geq0.
 \label{eq:clifford-transfer-ratio}
\end{align}
Thus the same \(M\) is available for iteration; see
\cref{cor:iterated-clifford-transfer}. In particular, \(a=1\)
preserves polynomial potentials and gives the sufficient bound
\(N\geq5\ell+3+(4\ell+2)M\).

The harmonic radial criterion supplies further target potentials.
Let \(f\) be homogeneous of degree \(e\) on \(\R^m\) and satisfy
the hypotheses of \cref{cor:harmonic-radial-existence}. Set
\(\sigma=(m-1-2e)/2>0\) and \(H=|y|^\sigma f\). Harmonicity gives
\[
 H\Delta_yH
 =\sigma(\sigma+2e+m-2)|y|^{2\sigma-2}f^2\geq0
 \qquad(y\ne0).
\]
Together with the corollary's phase sign and nonvanishing gradient,
this verifies all target hypotheses with \(Z=\{0\}\),
\(\ell=e+\sigma\), and \(M=0\). Thus
\cref{thm:clifford-area-transfer} applies whenever the source-dimension
bound \eqref{eq:clifford-transfer-bound} holds.
